(1) 两台理想50:50无损分束器与两臂组成Mach–Zehnder干涉仪。采用 \(U_{\rm BS}=2^{-1/2}\begin{pmatrix}1&i\\i&1\end{pmatrix}\)，输入复振幅为 \((E_0,0)^T\)，求第一次分束后两臂振幅。
(2) 两臂分别积累相位 \(\phi_1,\phi_2\)，求第二次分束之后的两个输出复振幅。
(3) 计算两输出端光强比例，证明和为一。
(4) 若一臂额外伸长光程 \(\Delta\ell\)，求两端强度关于 \(\Delta\ell\) 的周期，并说明零臂差极限。
